\section{区分求積法}\label{節:和の区分求積}

次の和は,\,$n$を大きくするとどうなるだろうか.
\[
 S_n=\sum_{k=1}^{n}\frac{n}{n^2+k^2}
\]
各項の分母は$n^2$以上なので,どの項も$\frac1n$以下であり,\,$n\to\infty$で$0$へ近づく.一方で,足す項の個数$n$は増えていく.どちらの影響が強いのかは,式を眺めるだけでは分からない.そこで実際に計算してみよう.
\begin{center}
\begin{tabular}{ccccccc}
\toprule
$n$ & $1$ & $2$ & $5$ & $10$ & $100$ & $1000$\\
\midrule
$S_n$ & $0.5$ & $0.65$ & $0.7337$ & $0.7600$ & $0.7829$ & $0.7851$\\
\bottomrule
\end{tabular}
\end{center}
$S_n$は$0$へは向かわず,\,$0.785$付近へ近づいているように見える.この値の正体を突き止めよう.

ここで1つの知識を思い出そう.定積分は,区間を細かく分けて作った長方形の面積の和の極限であった.そこで,和の極限を定積分に帰着させることを目標とする.

連続関数$f$について,区間$[a,b]$を$n$等分し,各小区間の右端での関数値と幅$\frac{b-a}{n}$を掛けて足すと,
\[
 \frac{b-a}{n}\sum_{k=1}^{n}
 f\left(a+\frac{k(b-a)}{n}\right)
\]
となる.\,$f$が0以上なら長方形の面積の和であり,負になる部分は符号を付けて足している.\,$n$を大きくすれば,この和は$\int_a^b f(x)\dd x$へ近づく.とくに$[0,1]$では
\begin{equation}\label{式:単位区間の区分求積}
 \lim_{n\to\infty}\frac1n\sum_{k=1}^{n}f\left(\frac{k}{n}\right)
 =\int_0^1 f(x)\dd x
\end{equation}
となる.この式を使って和の極限を求める方法を\textbf{区分求積法}という.

使うときには,\,$\frac kn$が横の位置,\,$\frac1n$が長方形の幅であることを確かめる.分母に$n$が現れるというだけでは足りない.和全体を,幅と高さの積の和へ書き直す必要がある.

\begin{bookpoint}
\textbf{POINT}\par
和から幅$\frac1n$をくくり出し,残りを$\frac kn$だけの式$f\left(\frac kn\right)$に書き直して,\,$\frac kn$が動く範囲で$f$を積分する.
\end{bookpoint}

\begin{example}{和から積分へ}
次の極限を求めよ.
\[
 \lim_{n\to\infty}\sum_{k=1}^{n}\frac{n}{n^2+k^2}.
\]
\end{example}

\paragraphbf{方針}
各項の分母から$n^2$をくくり,幅$\frac1n$と$\frac kn$の関数の積に書き直す.

\paragraphbf{解答}
$\frac{\pi}{4}$

\paragraphbf{解法}
各項の分母から$n^2$をくくると
\[
 \frac{n}{n^2+k^2}
 =\frac1n\cdot\cfrac{1}{1+\left(\cfrac{k}{n}\right)^2}
\]
となる.したがって,これは$f(x)=\frac1{1+x^2}$の右端を使う区分求積である.この関数は$[0,1]$で連続なので
\[
 \lim_{n\to\infty}\sum_{k=1}^{n}\frac{n}{n^2+k^2}
 =\int_0^1\frac{1}{1+x^2}\dd x
\]
を得る.右辺で$x=\tan t$と置くと,\,$x$が$0$から$1$まで動くとき$t$は$0$から$\frac\pi4$まで動き,\,$\dd x=(1+\tan^2t)\dd t$である.よって
\[
 \int_0^1\frac{1}{1+x^2}\dd x
 =\int_0^{\frac\pi4}\frac{1+\tan^2t}{1+\tan^2t}\dd t
 =\frac{\pi}{4}.
\]
表で見た$0.785$付近の値は,\,$\frac\pi4=0.78539\ldots$だったのだ.

\paragraphbf{積分区間の決定}
和に含まれる点の範囲が変わるなら,積分区間も確かめなければならない.例えば
\[
 \frac1n\sum_{k=1}^{2n}f\left(\frac{k}{n}\right)
\]
では,幅は$\frac1n$のままで,点は$\frac1n$から$2$まで並ぶ.これは$[0,2]$の区分求積である.項数が$2n$だからといって,幅を$\frac1{2n}$に変えてはいけない.

\begin{example}{積分区間の決定}
次の極限を求めよ.
\[
 \lim_{n\to\infty}\sum_{k=n+1}^{2n}\frac1k.
\]
\end{example}

\paragraphbf{方針}
各項を$\frac1n\cdot\cfrac1{\cfrac kn}$と書き,\,$k=n+1,\ldots,2n$に対応する$\frac kn$の範囲から積分区間を決める.

\paragraphbf{解答}
$\log2$

\paragraphbf{解法}
各項を
\[
 \frac1k=\frac1n\cdot\cfrac1{\cfrac{k}{n}}
\]
と書く.\,$k=n+1,\ldots,2n$に対応する点$\frac kn$は$1+\frac1n,\ldots,2$であり,\,$[1,2]$を$n$等分した小区間の右端である.\,$\frac1x$は$[1,2]$で連続なので
\[
 \lim_{n\to\infty}\sum_{k=n+1}^{2n}\frac1k
 =\int_1^2\frac1x\dd x
 =\log2.
\]
$\frac1x$は$x=0$で定義されないが,使った積分区間は$[1,2]$なので問題にならない.

\paragraphbf{端の項の扱い}
式\eqref{式:単位区間の区分求積}で,右端の代わりに左端や中点を使っても,同じ積分へ収束する.左端を使う和は,右端を使う和と端の1項ずつが違うだけである.実際,連続関数$f$について
\[
 \frac1n\sum_{k=0}^{n}f\left(\frac{k}{n}\right)
 -\frac1n\sum_{k=1}^{n}f\left(\frac{k}{n}\right)
 =\frac{f(0)}n
 \longrightarrow0
\]
である.このように,有限個の端の項を付けたり取り除いたりしても,その差が$0$へ近づくなら極限は変わらない.しかし,取り除いた項数が$n$とともに増える場合や,端の関数値が大きくなる場合には,この確認を省けない.区分求積の公式を使う範囲を先に決め,残りの差は別に評価すると良い.
